%% Generated by tools/gen_error_codes.py from docs/errors/ of the compiler;
%% edit the compiler's pages or tools/error_themes.txt, then rerun it.

\clearpage
\section{Operators and indexing}
\label{sec:codes:operators}

\begin{longtable}{@{}l p{0.78\linewidth} r@{}}
\toprule Code & Message & Page\\ \midrule \endhead
\hyperref[sec:codes:e4005]{E4005} & slice access, index overflow \errhole{} (len) \textless{}= \errhole{} (index) & \pageref{sec:codes:e4005}\\
\hyperref[sec:codes:e4040]{E4040} & range index cannot be decreasing for type \errhole{}, but goes from \errhole{} to \errhole{} & \pageref{sec:codes:e4040}\\
\hyperref[sec:codes:e4047]{E4047} & '\$' is usable only inside an index operation & \pageref{sec:codes:e4047}\\
\hyperref[sec:codes:e4142]{E4142} & index cannot be negative for type \errhole{}, but is \errhole{} & \pageref{sec:codes:e4142}\\
\hyperref[sec:codes:e4195]{E4195} & the index operator on \errhole{} with a dynamic operand of type \errhole{} is allowed only in the context of a copy statement & \pageref{sec:codes:e4195}\\
\hyperref[sec:codes:e4222]{E4222} & undefined operator \errhole{} for type \errhole{} and field \errhole{} & \pageref{sec:codes:e4222}\\
\hyperref[sec:codes:e4223]{E4223} & undefined module operator \errhole{} for \errhole{} and field \errhole{} & \pageref{sec:codes:e4223}\\
\hyperref[sec:codes:e4224]{E4224} & undefined operator \errhole{} for types \errhole{} and \errhole{} & \pageref{sec:codes:e4224}\\
\hyperref[sec:codes:e4229]{E4229} & operator \$ is not defined for type \errhole{} & \pageref{sec:codes:e4229}\\
\hyperref[sec:codes:e4231]{E4231} & the index operator is not defined for type \errhole{} and \errhole{} & \pageref{sec:codes:e4231}\\
\hyperref[sec:codes:e4237]{E4237} & undefined operator \errhole{} for type \errhole{} & \pageref{sec:codes:e4237}\\
\hyperref[sec:codes:e5011]{E5011} & slice access always overflows, \errhole{} (len) \errhole{} \errhole{} (index) & \pageref{sec:codes:e5011}\\
\bottomrule
\end{longtable}

\errorcode{E4005}{slice access, index overflow \errhole{} (len) \textless{}= \errhole{} (index)}
\label{sec:codes:e4005}

Raised during \textbf{validation}, from \texttt{ValidateErrorMessage::\allowbreak{}ARRAY\_\allowbreak{}OVERFLOW} (\textit{src/\allowbreak{}ymirc/\allowbreak{}semantic/\allowbreak{}validator/\allowbreak{}errors.yr}).

\subsubsection*{What it means}

A slice is indexed at a position the compiler could compute, and that position is outside the slice. Both the length and the index are in the message.

\subsubsection*{Example}

\textit{test\_\allowbreak{}resources/\allowbreak{}lit\_\allowbreak{}macro/\allowbreak{}test29.yr}:

%% check: skip
\begin{lstlisting}[style=coloredverbatimError]
/*
 * An indexed access of a repeated capture, that is out of the bounds of
 * what the call captured
 */

macro pick {
    self (v=(a=__operand)*) skips (" ") #{
        #(v[5]::a)
    }
}

fn foo ()-> i32 {
    pick#{ 1 2 }
}
\end{lstlisting}

\begin{lstlisting}[style=bashVerb, breaklines=true, escapechar=~]
Error[E4005] : slice access, index overflow 2 (len) <= 5 (index)
 --> :(1,8)
 1  ~\lstbox{\textSFxi{}}~     #(v[5]::a)
    ~\lstbox{\textSFv{}}~        ^
    ~\lstbox{\textSFii{}}~~\lstbox{\textSFx{}}~~\lstbox{\textSFx{}}~~\lstbox{\textSFx{}}~~\lstbox{\textSFx{}}~~\lstbox{\textSFx{}}~~\lstbox{\textSFx{}}~~\lstbox{\textSFx{}}~
\end{lstlisting}

\subsubsection*{How to fix it}

Correct the index, or check the length at run time before indexing when it is not known statically.

\errorcode{E4040}{range index cannot be decreasing for type \errhole{}, but goes from \errhole{} to \errhole{}}
\label{sec:codes:e4040}

Raised during \textbf{validation}, from \texttt{ValidateErrorMessage::\allowbreak{}DECREASING\_\allowbreak{}RANGE\_\allowbreak{}ACCESS} (\textit{src/\allowbreak{}ymirc/\allowbreak{}semantic/\allowbreak{}validator/\allowbreak{}errors.yr}).

\subsubsection*{What it means}

A range index goes from a higher bound to a lower one on a type whose indices only increase.

\subsubsection*{Example}

\textit{test\_\allowbreak{}resources/\allowbreak{}diagnostics/\allowbreak{}test26.yr}:

%% check: skip
\begin{lstlisting}[style=coloredverbatimError]
fn main () {
    let a = [1, 2, 3];
    let _ = a [2 .. 0];
}
\end{lstlisting}

\begin{lstlisting}[style=bashVerb, breaklines=true, escapechar=~]
Error[E4040] : range index cannot be decreasing for type [i32 ; 3us], but goes from 2 to 0
 --> test_resources/diagnostics/test26.yr:(3,15)
 3  ~\lstbox{\textSFxi{}}~     let _ = a [2 .. 0];
    ~\lstbox{\textSFv{}}~               ^  ^^
    ~\lstbox{\textSFii{}}~~\lstbox{\textSFx{}}~~\lstbox{\textSFx{}}~~\lstbox{\textSFx{}}~~\lstbox{\textSFx{}}~~\lstbox{\textSFx{}}~~\lstbox{\textSFx{}}~~\lstbox{\textSFx{}}~
\end{lstlisting}

\subsubsection*{How to fix it}

Swap the bounds. To walk the values backwards, reverse the result rather than the range.

\errorcode{E4047}{'\$' is usable only inside an index operation}
\label{sec:codes:e4047}

Raised during \textbf{validation}, from \texttt{ValidateErrorMessage::\allowbreak{}DOLLAR\_\allowbreak{}OUTSIDE\_\allowbreak{}CONTEXT} (\textit{src/\allowbreak{}ymirc/\allowbreak{}semantic/\allowbreak{}validator/\allowbreak{}errors.yr}).

\subsubsection*{What it means}

\tokennolst{\$}, which stands for the length of the value being indexed, was used outside an index operation, where it has nothing to refer to.

\subsubsection*{Example}

\textit{test\_\allowbreak{}resources/\allowbreak{}diagnostics/\allowbreak{}test29.yr}:

%% check: skip
\begin{lstlisting}[style=coloredverbatimError]
fn main () {
    let _ = $;
}
\end{lstlisting}

\begin{lstlisting}[style=bashVerb, breaklines=true, escapechar=~]
Error[E4047] : '$' is usable only inside an index operation
 --> test_resources/diagnostics/test29.yr:(2,13)
 2  ~\lstbox{\textSFxi{}}~     let _ = $;
    ~\lstbox{\textSFv{}}~             ^
    ~\lstbox{\textSFii{}}~~\lstbox{\textSFx{}}~~\lstbox{\textSFx{}}~~\lstbox{\textSFx{}}~~\lstbox{\textSFx{}}~~\lstbox{\textSFx{}}~~\lstbox{\textSFx{}}~~\lstbox{\textSFx{}}~
\end{lstlisting}

\subsubsection*{How to fix it}

Use the length explicitly (\tokennolst{x.len}) outside an index.

\errorcode{E4142}{index cannot be negative for type \errhole{}, but is \errhole{}}
\label{sec:codes:e4142}

Raised during \textbf{validation}, from \texttt{ValidateErrorMessage::\allowbreak{}NEGATIVE\_\allowbreak{}INT\_\allowbreak{}INDEX} (\textit{src/\allowbreak{}ymirc/\allowbreak{}semantic/\allowbreak{}validator/\allowbreak{}errors.yr}).

\subsubsection*{What it means}

An index the compiler computed is negative, on a type whose indices start at zero.

\subsubsection*{Example}

\textit{test\_\allowbreak{}resources/\allowbreak{}diagnostics/\allowbreak{}test56.yr}:

%% check: skip
\begin{lstlisting}[style=coloredverbatimError]
fn main () {
    let a = [1, 2, 3];
    let _ = a [-1];
}
\end{lstlisting}

\begin{lstlisting}[style=bashVerb, breaklines=true, escapechar=~]
Error[E4142] : index cannot be negative for type [i32 ; 3us], but is -1
 --> test_resources/diagnostics/test56.yr:(3,15)
 3  ~\lstbox{\textSFxi{}}~     let _ = a [-1];
    ~\lstbox{\textSFv{}}~               ^^
    ~\lstbox{\textSFii{}}~~\lstbox{\textSFx{}}~~\lstbox{\textSFx{}}~~\lstbox{\textSFx{}}~~\lstbox{\textSFx{}}~~\lstbox{\textSFx{}}~~\lstbox{\textSFx{}}~~\lstbox{\textSFx{}}~
\end{lstlisting}

\subsubsection*{How to fix it}

Correct the index. To count from the end, use \tokennolst{\$} inside the index operation.

\errorcode{E4195}{the index operator on \errhole{} with a dynamic operand of type \errhole{} is allowed only in the context of a copy statement}
\label{sec:codes:e4195}

Raised during \textbf{validation}, from \texttt{ValidateErrorMessage::\allowbreak{}RANGE\_\allowbreak{}ON\_\allowbreak{}ARRAY\_\allowbreak{}NO\_\allowbreak{}COPY} (\textit{src/\allowbreak{}ymirc/\allowbreak{}semantic/\allowbreak{}validator/\allowbreak{}errors.yr}).

\subsubsection*{What it means}

A range index whose bounds are not known at compile time was applied outside a copy. Such an index produces a new value rather than a view, which only the copy context accepts.

\subsubsection*{Example}

\textit{test\_\allowbreak{}resources/\allowbreak{}diagnostics/\allowbreak{}test122.yr}:

%% check: skip
\begin{lstlisting}[style=coloredverbatimError]
// a range index with a runtime bound on a static array, outside a copy
fn main () {
    let a = [1, 2, 3];
    let mut i = 0us;
    i = 1us;
    let _ = a [i .. 2us];
}
\end{lstlisting}

\begin{lstlisting}[style=bashVerb, breaklines=true, escapechar=~]
Error[E4195] : the index operator on [i32 ; 3us] with a dynamic operand of type (..usize) is allowed only in the context of a copy statement
 --> test_resources/diagnostics/test122.yr:(6,15)
 6  ~\lstbox{\textSFxi{}}~     let _ = a [i .. 2us];
    ~\lstbox{\textSFv{}}~               ^  error: value of type (..usize) is needed but unknown at compilation time
    ~\lstbox{\textSFxi{}}~
    = when validating test122::main -- (test_resources/diagnostics/test122.yr:(2,4))
    ~\lstbox{\textSFii{}}~~\lstbox{\textSFx{}}~~\lstbox{\textSFx{}}~~\lstbox{\textSFx{}}~~\lstbox{\textSFx{}}~~\lstbox{\textSFx{}}~~\lstbox{\textSFx{}}~~\lstbox{\textSFx{}}~
\end{lstlisting}

\subsubsection*{How to fix it}

Enclose the index in a \tokennolst{copy}, or use bounds the compiler can compute.

\errorcode{E4222}{undefined operator \errhole{} for type \errhole{} and field \errhole{}}
\label{sec:codes:e4222}

Raised during \textbf{validation}, from \texttt{ValidateErrorMessage::\allowbreak{}UNDEFINED\_\allowbreak{}BIN\_\allowbreak{}ACC\_\allowbreak{}OP} (\textit{src/\allowbreak{}ymirc/\allowbreak{}semantic/\allowbreak{}validator/\allowbreak{}errors.yr}).

\subsubsection*{What it means}

A binary operator is not defined between the type and the field named.

\subsubsection*{Example}

\textit{test\_\allowbreak{}resources/\allowbreak{}regressions/\allowbreak{}float\_\allowbreak{}properties/\allowbreak{}test2.yr}:

%% check: skip
\begin{lstlisting}[style=coloredverbatimError]
in test2;

fn main () {

    f32::foo;
    f64::bar;
    f80::baz;
    fsize::bad;


}
\end{lstlisting}

\begin{lstlisting}[style=bashVerb, breaklines=true, escapechar=~]
Error[E4222] : undefined operator :: for type f32 and field foo
 --> test_resources/regressions/float_properties/test2.yr:(5,8)
 5  ~\lstbox{\textSFxi{}}~     f32::foo;
    ~\lstbox{\textSFv{}}~        ^^
    ~\lstbox{\textSFii{}}~~\lstbox{\textSFx{}}~~\lstbox{\textSFx{}}~~\lstbox{\textSFx{}}~~\lstbox{\textSFx{}}~~\lstbox{\textSFx{}}~~\lstbox{\textSFx{}}~~\lstbox{\textSFx{}}~
\end{lstlisting}

\subsubsection*{Notes}

When the field is \tokennolst{typeinfo}, the reason is reported under this diagnostic, as a note carrying no code of its own:

\begin{itemize}
\item \texttt{temporary type \textless{}\textgreater{} has no type information} (formerly E4219): \tokennolst{typeinfo} was requested on a type that has none, such as a trait.
\end{itemize}

\subsubsection*{How to fix it}

Check the field: a field method may be needed, or the operator may have to be defined on the field type.

\errorcode{E4223}{undefined module operator \errhole{} for \errhole{} and field \errhole{}}
\label{sec:codes:e4223}

Raised during \textbf{validation}, from \texttt{ValidateErrorMessage::\allowbreak{}UNDEFINED\_\allowbreak{}BIN\_\allowbreak{}MOD\_\allowbreak{}OP} (\textit{src/\allowbreak{}ymirc/\allowbreak{}semantic/\allowbreak{}validator/\allowbreak{}errors.yr}).

\subsubsection*{What it means}

A module-level binary operator is not defined for the operands given.

\subsubsection*{Example}

\textit{test\_\allowbreak{}resources/\allowbreak{}diagnostics/\allowbreak{}test62.yr}:

%% check: skip
\begin{lstlisting}[style=coloredverbatimError]
class Foo {
    pub self () {}
}

fn main () { let _ = copy Foo::named (); }
\end{lstlisting}

\begin{lstlisting}[style=bashVerb, breaklines=true, escapechar=~]
Error[E4223] : undefined module operator :: for test62::Foo::self and field named
 --> test_resources/diagnostics/test62.yr:(5,30)
 5  ~\lstbox{\textSFxi{}}~ fn main () { let _ = copy Foo::named (); }
    ~\lstbox{\textSFv{}}~                              ^^
    ~\lstbox{\textSFii{}}~~\lstbox{\textSFx{}}~~\lstbox{\textSFx{}}~~\lstbox{\textSFx{}}~~\lstbox{\textSFx{}}~~\lstbox{\textSFx{}}~~\lstbox{\textSFx{}}~~\lstbox{\textSFx{}}~
\end{lstlisting}

\subsubsection*{How to fix it}

Define the operator, or use one that exists for these types.

\errorcode{E4224}{undefined operator \errhole{} for types \errhole{} and \errhole{}}
\label{sec:codes:e4224}

Raised during \textbf{validation}, from \texttt{ValidateErrorMessage::\allowbreak{}UNDEFINED\_\allowbreak{}BIN\_\allowbreak{}OP} (\textit{src/\allowbreak{}ymirc/\allowbreak{}semantic/\allowbreak{}validator/\allowbreak{}errors.yr}).

\subsubsection*{What it means}

A binary operator is not defined for the two operand types. Both are in the message.

\subsubsection*{Example}

\textit{test\_\allowbreak{}resources/\allowbreak{}lit\_\allowbreak{}class/\allowbreak{}operators/\allowbreak{}test6.yr}:

%% check: skip
\begin{lstlisting}[style=coloredverbatimError]
class A {
    pub self () {}
}

fn main () {
    let a = copy A ();
    a * a;

}
\end{lstlisting}

\begin{lstlisting}[style=bashVerb, breaklines=true, escapechar=~]
Error[E4224] : undefined operator * for types &(test6::A) and &(test6::A)
 --> test_resources/lit_class/operators/test6.yr:(7,7)
 7  ~\lstbox{\textSFxi{}}~     a * a;
    ~\lstbox{\textSFv{}}~       ^  error: class &(test6::A) has no method named opBinary
    ~\lstbox{\textSFxi{}}~       ~\lstbox{\textSFii{}}~~\lstbox{\textSFx{}}~ error: class &(test6::A) has no method named opBinaryRight
    ~\lstbox{\textSFxi{}}~
    = when validating test6::main -- (test_resources/lit_class/operators/test6.yr:(5,4))
    ~\lstbox{\textSFii{}}~~\lstbox{\textSFx{}}~~\lstbox{\textSFx{}}~~\lstbox{\textSFx{}}~~\lstbox{\textSFx{}}~~\lstbox{\textSFx{}}~~\lstbox{\textSFx{}}~~\lstbox{\textSFx{}}~
\end{lstlisting}

\subsubsection*{How to fix it}

Convert one operand, or define the operator on the type. When both types look the same, the difference is in their mutability, which is part of the type.

\errorcode{E4229}{operator \$ is not defined for type \errhole{}}
\label{sec:codes:e4229}

Raised during \textbf{validation}, from \texttt{ValidateErrorMessage::\allowbreak{}UNDEFINED\_\allowbreak{}DOLLAR\_\allowbreak{}OP} (\textit{src/\allowbreak{}ymirc/\allowbreak{}semantic/\allowbreak{}validator/\allowbreak{}errors.yr}).

\subsubsection*{What it means}

\tokennolst{\$} was used on a type that does not define a length.

\subsubsection*{Example}

\textit{test\_\allowbreak{}resources/\allowbreak{}diagnostics/\allowbreak{}test74.yr}:

%% check: skip
\begin{lstlisting}[style=coloredverbatimError]
class Foo { pub self () {} }

fn main () {
    let f = copy Foo ();
    let _ = f [$];
}
\end{lstlisting}

\begin{lstlisting}[style=bashVerb, breaklines=true, escapechar=~]
Error[E4229] : operator $ is not defined for type &(test74::Foo)
 --> test_resources/diagnostics/test74.yr:(5,16)
 5  ~\lstbox{\textSFxi{}}~     let _ = f [$];
    ~\lstbox{\textSFv{}}~                ^  error: class &(test74::Foo) has no method named opDollar
    ~\lstbox{\textSFxi{}}~
    = when validating test74::main -- (test_resources/diagnostics/test74.yr:(3,4))
    ~\lstbox{\textSFii{}}~~\lstbox{\textSFx{}}~~\lstbox{\textSFx{}}~~\lstbox{\textSFx{}}~~\lstbox{\textSFx{}}~~\lstbox{\textSFx{}}~~\lstbox{\textSFx{}}~~\lstbox{\textSFx{}}~
\end{lstlisting}

\subsubsection*{How to fix it}

Use a type with a length, or write the bound explicitly.

\errorcode{E4231}{the index operator is not defined for type \errhole{} and \errhole{}}
\label{sec:codes:e4231}

Raised during \textbf{validation}, from \texttt{ValidateErrorMessage::\allowbreak{}UNDEFINED\_\allowbreak{}INDEX\_\allowbreak{}OP} (\textit{src/\allowbreak{}ymirc/\allowbreak{}semantic/\allowbreak{}validator/\allowbreak{}errors.yr}).

\subsubsection*{What it means}

The index operator is not defined between the indexed type and the index type. Both are in the message.

\subsubsection*{Example}

\textit{test\_\allowbreak{}resources/\allowbreak{}aka/\allowbreak{}test2.yr}:

%% check: skip
\begin{lstlisting}[style=coloredverbatimError]
def A : dmut [i32];

fn main ()
{
    let dmut x : A = copy [1, 2, 3];
    x [0] = 1;
}
\end{lstlisting}

\begin{lstlisting}[style=bashVerb, breaklines=true, escapechar=~]
Error[E4231] : the index operator is not defined for type error and {i32}
 --> test_resources/aka/test2.yr:(6,7)
 6  ~\lstbox{\textSFxi{}}~     x [0] = 1;
    ~\lstbox{\textSFv{}}~       ^
    ~\lstbox{\textSFii{}}~~\lstbox{\textSFx{}}~~\lstbox{\textSFx{}}~~\lstbox{\textSFx{}}~~\lstbox{\textSFx{}}~~\lstbox{\textSFx{}}~~\lstbox{\textSFx{}}~~\lstbox{\textSFx{}}~
\end{lstlisting}

\subsubsection*{How to fix it}

Index with the type the operation expects, or define the index operator on the type.

\errorcode{E4237}{undefined operator \errhole{} for type \errhole{}}
\label{sec:codes:e4237}

Raised during \textbf{validation}, from \texttt{ValidateErrorMessage::\allowbreak{}UNDEFINED\_\allowbreak{}UN\_\allowbreak{}OP} (\textit{src/\allowbreak{}ymirc/\allowbreak{}semantic/\allowbreak{}validator/\allowbreak{}errors.yr}).

\subsubsection*{What it means}

A unary operator is not defined for the operand type.

\subsubsection*{Example}

\textit{test\_\allowbreak{}resources/\allowbreak{}char\_\allowbreak{}type/\allowbreak{}test4.yr}:

%% check: skip
\begin{lstlisting}[style=coloredverbatimError]
in test4;


fn main () {
    let a = 'r'c8;
    let b = 'r'c16;
    let c = 'r'c32;


    *a;
    !a;
    -a;
    &a;

    *b;
    !b;
    -b;
    &b;

    *c;
    !c;
    -c;
    &c;
}
\end{lstlisting}

\begin{lstlisting}[style=bashVerb, breaklines=true, escapechar=~]
Error[E4237] : undefined operator * for type c8
 --> test_resources/char_type/test4.yr:(10,5)
10  ~\lstbox{\textSFxi{}}~     *a;
    ~\lstbox{\textSFv{}}~     ^
    ~\lstbox{\textSFii{}}~~\lstbox{\textSFx{}}~~\lstbox{\textSFx{}}~~\lstbox{\textSFx{}}~~\lstbox{\textSFx{}}~~\lstbox{\textSFx{}}~~\lstbox{\textSFx{}}~~\lstbox{\textSFx{}}~
\end{lstlisting}

\subsubsection*{How to fix it}

Define the operator on the type, or convert the operand.

\errorcode{E5011}{slice access always overflows, \errhole{} (len) \errhole{} \errhole{} (index)}
\label{sec:codes:e5011}

Raised during \textbf{lowering (YIL)}, from \tokennolst{OptimizerErrorMessage::ALWAYS\_OVERFLOW} (\textit{src/\allowbreak{}ymirc/\allowbreak{}lint/\allowbreak{}optimizer/\allowbreak{}errors.yr}), by the bounds check pass of the optimizer.

\subsubsection*{What it means}

An access to a slice or an array is out of its bounds on every execution that reaches it: the check guarding it is always taken, and the program panics there.

The validator already rejects an index it can compute on a length it can compute (\hyperref[sec:codes:e4005]{E4005}). This one is found later, on what the control flow of the function says: the length of a mutable slice after the last assignment to it, or the length a slice was given on every path leading to the access. The optimizer only runs from \tokennolst{-O1}, so the diagnostic is never raised at \tokennolst{-O0}.

It is a warning: the access may sit on a path that never runs. Like any warning, it only lets the compilation carry on in debug mode. The check and the panic it guards are kept, so a program that reaches the access still panics there.

\subsubsection*{Example}

\textit{test\_\allowbreak{}resources/\allowbreak{}warnings/\allowbreak{}test31.yr}:

%% check: skip
\begin{lstlisting}[style=coloredverbatimError]
// an access the optimizer proves out of bounds warns, and panics as it did
fn foo ()-> i32 {
    let dmut a = copy [1, 2, 3];

    a [10]
}
\end{lstlisting}

\begin{lstlisting}[style=bashVerb, breaklines=true, escapechar=~]
Warning[E5011] : slice access always overflows, 3 (len) <= 10 (index)
 --> test_resources/warnings/test31.yr:(5,7)
 5  ~\lstbox{\textSFxi{}}~     a [10]
    ~\lstbox{\textSFv{}}~       ^--
    ~\lstbox{\textSFii{}}~~\lstbox{\textSFx{}}~~\lstbox{\textSFx{}}~~\lstbox{\textSFx{}}~~\lstbox{\textSFx{}}~~\lstbox{\textSFx{}}~~\lstbox{\textSFx{}}~~\lstbox{\textSFx{}}~
\end{lstlisting}

\subsubsection*{How to fix it}

Correct the index, or the length the slice is given before it is read.
